\chapter{Исследовательская часть}
%<цель исследования, на какой машине делали (указать ЦПУ, ОЗУ, ОС), желательно написать, как исследовали и при каких условиях; что получили в результате (таблицы + графики)>

\section{Замеры процессорного времени}

Замеры проводились на ноутбуке Redmibook Pro 14 2022:
\begin{itemize}
	\item процессор Intel i7-12650H (архитектура x86-64, 10 ядер по 2.3 Ггц);
	\item оперативная память 16 Гб (DDR4, 5200 МГц);
	\item операционная система Windows 11 Home 24H2.
\end{itemize}
Для замеров использовалась функция clock() из библиотеки ctime~\cite[с.~288]{lit4}. 

% выравнивание по точке
\begin{table}[!htbp]
	\caption{Некоторые замеры процессорного времени для квадратных матриц.}
	\centering
	\begin{tabularx}{\textwidth}{|>{\centering\arraybackslash}X|*{3}{>{\hspace{2em}}X}}%S[table-format=4.5]|}}
		\hline
		\textbf{Размер} & \textbf{Стандартный} & \textbf{Виноград} & \textbf{Виноград с оптимизацией} \\\hline
		2  & 0.0005 & 0.0005 & 0.0005 \\\hline
		3  & 0.0005 & 0.0005 & 0.0010 \\\hline
		9  & 0.0120 & 0.0085 & 0.0055 \\\hline
		10 & 0.0105 & 0.0105 & 0.0090 \\\hline
		49 & 1.0710 & 1.0265 & 0.6920 \\\hline
		50 & 1.0885 & 1.1060 & 0.7635 \\\hline
	\end{tabularx}
\end{table}

% график должен быть читаем в чб
% графики иллюстрируют

% таблица + график из txt/csv

\begin{figure}[!htbp]
	\centering
	\begin{tikzpicture}
		\selectcolormodel{gray}
		\begin{axis}[
			axis lines = left, 
			xlabel = Размер, 
			ylabel = Время, 
			grid=both,
			minor grid style={gray!25}, 
			major grid style={gray!50},
			legend style={font=\scriptsize, legend pos=north west}
		]
			\addplot+[smooth, mark=none] table [x index=0, y index=1] {code/proctimeres.txt};
			\addlegendentry{Стандартный}
			\addplot+[smooth, mark=none, dashed] table [x index=0, y index=2] {code/proctimeres.txt};
			\addlegendentry{Виноград}
			\addplot+[smooth, mark=none, thick] table [x index=0, y index=3] {code/proctimeres.txt};
			\addlegendentry{Виноград с оптимизацией}
		\end{axis}
	\end{tikzpicture}
	\caption{График зависимости процессорного времени от размера матрицы}
\end{figure}

\section*{Вывод}
%<что сделали и получили в результате, кратко>

В данном разделе были проведены замеры процессорного времени для каждой реализации из предыдущего раздела. Полученные на основе этих замеров графики иллюстрируют рассчитанные ранее трудоёмкости рассматриваемых в работе алгоритмов.

\clearpage
